\documentclass[10pt,a4paper]{amsart}
\usepackage[margin=19mm]{geometry}
\usepackage{amsmath,amssymb,mathtools}
\usepackage[scr=rsfs,cal=euler]{mathalfa}
\usepackage{xfrac}
\usepackage[unicode,hidelinks,bookmarks=false]{hyperref}
\newcommand{\ie}{i.e.\ }
\newcommand{\cf}{cf.\ }
\newcommand{\resp}{respectively\ }
\newcommand{\Subsection}{\subsection}
\providecommand{\nobreakdash}{\nobreak\hbox{-}}
\setlength{\emergencystretch}{2em}
\multlinegap0pt
\usepackage{bm}
\usepackage{bbm}
\usepackage{dsfont}
\usepackage{xspace}
\usepackage{cancel}
\usepackage{mathtools}
\usepackage{amsthm}
\makeatletter
\@ifundefined{theorem}{\newtheorem{theorem}{Theorem}}{}
\makeatother
\title[Cosets of normal subgroups and union of two conjugacy classe]{Cosets of normal subgroups and union of two conjugacy classes}
\author{Antonio Beltr\'an}
\date{Source: arXiv:2508.06367v1 (CC BY 4.0)}
\begin{document}
\maketitle

\noindent\emph{Source and licence.} Cosets of normal subgroups and union of two conjugacy classes. Version arXiv:2508.06367v1 (CC BY 4.0), distributed under the Creative Commons Attribution 4.0 International licence, as recorded in the arXiv licence metadata for this exact version and shown on the abstract page https://arxiv.org/abs/2508.06367v1. Full rights notice: licenses/arxiv-2508.06367-CC-BY-4.0.txt.\par\smallskip

\noindent\textit{Editorial scope.} Counted unit: the Theorem whose source label is \texttt{2} in the article (source0.tex lines 141--153, source label \texttt{2}), delivered together with the article's own complete original proof of it (source0.tex lines 155--233, one \texttt{proof} environment of 79 lines). No mathematics is written, repaired or re-derived: statement and proof are the article's own text line for line. Required context retained uncounted: none. Every definition, display and label of those lines is retained unchanged. Omitted from this package and not redistributed: 1--140, 154--154, 234--641. Nothing inside a retained excerpt is omitted or altered: every retained body is delivered exactly as the article's lines are written, so no omission receipt of any kind accompanies this package. Citations: the retained excerpts carry no citation command at all, so no bibliography entry of the article is transcribed and no reference is redistributed. Reference adaptation: none was needed and none was made; every reference inside the delivered unit resolves inside it, so no unresolved reference or citation is produced. Printed slips kept literal: no printed word, symbol, hypothesis or number of the counted statement or of its proof was altered. Layout and class: the article's own class is \texttt{amsart} with packages including amssymb, enumerate, amsmath, xcolor; the wrapper is the lane's pinned \texttt{amsart} editorial wrapper at 10pt on A4, which restates the theorem-like environments the retained text needs and adds the article's own macro definitions that the retained text needs (none), with extra wrapper packages (bm, bbm, dsfont, xspace, cancel, mathtools). The delivered numbering is the unit's own. Assets: no retained span contains a figure environment, a graphics-inclusion command, a PDF-page inclusion, a table or a TikZ picture; no image file is copied into this package, the package declares no asset dependency and no image file is redistributed with it. Licence: version arXiv:2508.06367v1 of this article is distributed under the Creative Commons Attribution 4.0 International licence, as recorded in the arXiv licence metadata for this exact version and shown on the abstract page https://arxiv.org/abs/2508.06367v1. A full rights notice is delivered at licenses/arxiv-2508.06367-CC-BY-4.0.txt. Characters: every retained excerpt is pure ASCII, so the delivered engine, XeLaTeX with its default Latin Modern text font, reports no missing character.
 \begin{theorem}\label{2} Let $G$ be a finite group and $N$ a normal subgroup of $G$. Let $K =x^G$ and $D=d^G$ be distinct conjugacy classes of $x,d \in G$. Then the following conditions are equivalent:

 \begin{itemize}
 \item[(a)] $Nx \subseteq K \cup D$, but $Nx \not\subseteq K$.
 \item[(b)] $NK=ND=K\cup D$.
 \item[(c)] Every  $\chi \in {\rm Irr}(G/N)$ satisfies $\chi(x)=\chi(d)$; and  every $\chi\in {\rm Irr}(G|N)$ satisfies
 $$|K|\chi(x) +|D|\chi(d)=0,$$
and, in addition, $|{\bf C}_{G/N}(Nx)|= |G|/(|K|+|D|)$.
 \end{itemize}



\end{theorem}

\begin{proof}  The  equivalence between (a) and (b) is elementary. Assume (a), so there exist $n\in N$ and $d^g\in D$ for some $g\in G$, such that $nx=d^g$. Hence $Nx=Nd^g=(Nd)^g$. It clearly follows that $K\cup D\subseteq NK=ND\subseteq K\cup D$, so the equality holds. Conversely, if we assume (b), we trivially have $Nx\subseteq K\cup D$. Now the fact that  $Nx\subseteq K$ yields to $NK=K$, against  our assumption. Therefore, $Nx\not\subseteq K$, as required.


\medskip
We prove now that (a) and (b) imply (c). As we have said above, the hypotheses imply that $Nx=(Nd)^g$ for some $g\in G$, so
it is evident that $\chi(x)=\chi(d)$ for every $\chi\in {\rm Irr}(G/N)$.
Suppose now that $\chi\in{\rm Irr}(G|N)$ and let $\mathfrak{X}$ be a representation of $G$ that affords $\chi$.   Also,  we know that $\mathfrak{X}$ can be linearly extended to $\mathbb{C}[G]$.   Since $N$ is a disjoint union of conjugacy classes of $G$, then $$\widehat{ N} = \sum _{n \in N} n\in {\bf Z}(\mathbb{C}[G])$$ so, by Schur's Lemma, $\mathfrak{X}(\widehat{N})$ is a scalar matrix. On the other hand, the trace of $\mathfrak{X}(\widehat{N})$ is $$\sum_{n \in N}\chi(n)=|N|[\chi_N, 1_N] =0,$$ so $\mathfrak{X}(\widehat{N}) = O$, where $O$ denotes the zero matrix. Then $\mathfrak{X}(\widehat{Ng})=\mathfrak{X}(\widehat{N})\mathfrak{X}(g)=O$, for every $g\in G$. Now, by hypothesis, we can write $\widehat{K}\widehat{N} = m_1\widehat{K}+m_2\widehat{D}$ with $m_1, m_2>0$, and we know that $\mathfrak{X}(\widehat{N}\widehat{K})=\sum_{x^g \in K} \mathfrak{X}(\widehat{Nx^g})=O$.
Hence, by taking traces, we have
\begin{equation}
0={\rm Trace}(\mathfrak{X}(\widehat{N}\widehat{K}))=m_1 {\rm Trace}(\widehat{K}) + m_2 {\rm Trace}(\widehat{D})=m_1|K| \chi(x) + m_2|D| \chi(d). \label{eqn1}
\end{equation}

We will prove that $m_1=m_2$, and thus,  we will obtain $|K| \chi(x) +|D| \chi(d)=0$.
The second orthogonality relation together with the fact that  $\chi(x)=\chi(d)$ for every $\chi\in {\rm Irr}(G/N)$  give rise to the equations:

\begin{equation}
|{\bf C}_G(x)|=\sum_{\chi\in {\rm Irr}(G/N)} |\chi(x)|^2
+ \sum_{\chi\in {\rm Irr}(G|N)}
 \chi(x)\overline{\chi(x)}=|{\bf C}_{G/N}(Nx)|+ \sum_{\chi\in {\rm Irr}(G|N)}|\chi(x)|^2,  \label{eqn2}
 \end{equation}

\begin{equation} 0=\sum_{\chi\in {\rm Irr}(G/N)} |\chi(x)|^2
+ \sum_{\chi\in {\rm Irr}(G|N)}
 \chi(x)\overline{\chi(d)} =
|{\bf C}_{G/N}(Nx)| + \sum_{\chi\in {\rm Irr}(G|N)} \chi(x)\overline{\chi(d)}. \label{eqn3}
 \end{equation}
By applying Eq. (\ref{eqn1}) to Eq. (\ref{eqn3}), we get




\begin{equation}\label{eqn4}
0=|{\bf C}_{G/N}(Nx)| - \frac{m_1|K|}{m_2|D|}\sum_{\chi\in {\rm Irr}(G|N)}
 |\chi(x)|^2.
 \end{equation}

\noindent
On the other hand, the  equality $NK=K \cup D$ allows to write
\begin{equation}
Nx_1\cup \cdots \cup Nx_t= K\cup D \label{eqn5}
\end{equation}
where $Nx_i$ for $i=1,\ldots t$ are the distinct conjugates of $Nx$  in $G/N$, that is, $t=|G/N:{\bf C}_{G/N}(Nx)|$. Taking cardinalities in Eq. (\ref{eqn5}), we get $|N|t= |K|+ |D|$, and hence
$$|{\bf C}_{G/N}(Nx)|=  \frac{|G|}{|K|+ |D|}.$$
Thus, the last equality of (c) is proved.
We now substitute this value, together with the equality $|{\bf C}_G(x)|=|G|/|K|$, into  Eqs. (\ref{eqn2}) and (\ref{eqn4}),  and this gives
$$0= \frac{|G|}{|K|+ |D|}-\frac{m_1|K|}{m_2|D|}\left(\frac{|G|}{|K|}-\frac{|G|}{|K|+ |D|}\right).$$

\noindent
It follows that $m_1=m_2$, as required.

\medskip
In order to prove that (c) implies (b), we show first that $KC\subseteq K\cup D$ for every conjugacy class $C$ of $G$ contained in $N$. Write
 \begin{equation}\label{eqn6}
 \widehat{K}\widehat{C}=\alpha_K\widehat{K}+\alpha_D\widehat{D} + \sum_{i=1}^n \alpha_i \widehat{C_i},
 \end{equation}
 $C_1, \ldots , C_n$ being  all the conjugacy classes of $G$ (including the trivial class) distinct from $K$ and $D$, and $\alpha_K,\alpha_D, \alpha_i$ are non-negative integers. We compute  $\alpha_K|K| + \alpha_D|D|$.

 $$ \alpha_K|K| + \alpha_D|D|=\frac{|K|^2|C|}{|G|}\sum_{\chi\in {\rm Irr}(G)}\frac{\chi(x)\chi(c)\chi(x^{-1})}{\chi(1)}+
 \frac{|D||K||C|}{|G|}\sum_{\chi\in {\rm Irr}(G)}\frac{\chi(x)\chi(c)\chi(d^{-1})}{\chi(1)}=$$
 $$=\frac{|K|^2|C|}{|G|}\sum_{\chi\in {\rm Irr}(G)}\frac{|\chi(x)|^2\chi(c)}{\chi(1)}+
 \frac{|D||K||C|}{|G|}\sum_{\chi\in {\rm Irr}(G)}\frac{\chi(x)\chi(c)\overline{\chi(d})}{\chi(1)}= $$
 $$=\frac{|K|^2|C|}{|G|}(\sum_{\chi\in {\rm Irr}(G/N)}|\chi(x)|^2 + \sum_{\chi\in {\rm Irr}(G|N)}\frac{|\chi(x)|^2\chi(c)}{\chi(1)}) +$$
 $$+  \frac{|D||K||C|}{|G|}(\sum_{\chi\in {\rm Irr}(G/N)} \chi(x)\overline{\chi(d)} + \sum_{\chi\in {\rm Irr}(G|N)}\frac{\chi(x)\chi(c)\overline{\chi(d})}{\chi(1)}).$$

\noindent
By employing the  three equations appearing in statement (c), we have

  $$ \alpha_K|K| + \alpha_D|D|= (\frac{|K|^2|C|}{|G|} +
  \frac{|D||K||C|}{|G|})|{\bf C}_{G/N}(Nx)| +$$
 $$+ \frac{|K|^2|C|}{|G|}\sum_{\chi\in {\rm Irr}(G|N)}\frac{|\chi(x)|^2\chi(c)}{\chi(1)} +\frac{|D||K||C|}{|G|}
  \sum_{\chi\in {\rm Irr}(G|N)}(-\frac{|K|}{|D|})\frac{|\chi(x)|^2\chi(c)}{\chi(1)}=$$
$$=( \frac{|K|^2|C|}{|G|} +
  \frac{|D||K||C|}{|G|})\frac{|G|}{|K|+|D|}=|K||C|.$$
  
 On the other hand,  by counting elements in Eq. (\ref{eqn6}), we obtain
 $$|K||C|=\alpha_K|K|+ \alpha_D|D| + \sum_{i=1}^n\alpha_i|C_i|.$$
By joining the above results, we deduce that $\alpha_i=0$ for  $i=1,\ldots , n$. This means that $KC\subseteq K\cup D$, as required. We remark that $KC=K$ is feasible when $\alpha_D=0$ and $\alpha_K=|C|$ (for instance, when $C=1$),  and $KC=D$ can also occur (when $\alpha_K=0$ and $\alpha_D=|K||C|/|D|)$.
Accordingly, we conclude that $KN\subseteq K \cup D$. Finally, the fact that $\chi(x)=\chi(d)$ for every $\chi\in {\rm Irr}(G/N)$  is equivalent (because ${\rm Irr}(G)$ is a basis of the vector space of class functions of $G$) to assert that $Nx$ and $Nd$ are conjugate, and consequently, $NK=ND$. This obviously implies that $K\cup D\subseteq KN$, so the equality of (b) holds.
\end{proof}

\end{document}
